
Threshold signature schemes are a cryptographic primitive to facilitate joint
ownership over a private key by a set of participants,
such that a threshold number of participants must cooperate to issue a
signature that can be verified by a single public key.
Threshold signatures are useful
across a range of settings that require a distributed
root of trust among a set of equally trusted parties.

Similarly to signing operations in a single-party setting, some implementations
of threshold signature schemes require performing signing operations at
scale and under heavy load.
For example, threshold signatures can be used by a set of signers
to authenticate financial transactions in
cryptocurrencies~\cite{securing-bitcoin},
or to sign a
network consensus produced by a set of trusted authorities~\cite{pir-tor}.
In both of these examples, as the number of signing parties or signing
operations increases, the number of communication rounds between participants required
to produce the joint signature becomes a performance bottleneck, in addition to
the increased load experienced by each signer.
This problem is further exacerbated
when signers utilize network-limited devices or unreliable networks for
transmission, or protocols that wish to allow signers to participate in signing
operations asynchronously.
As such, optimizing the
network overhead of signing operations is highly beneficial to real-world
applications of threshold signatures.

Today in the literature, the best threshold signature schemes are those that
rely on pairing-based
cryptography~\cite{Boneh2004,BDDMNG:CompactMultiSignatures}, and can perform
signing operations in a single round among participants. However, relying on
pairing-based signature schemes is undesirable for some implementations in
practice, such as those that do not wish to introduce a new
cryptographic assumption, or that wish to maintain backwards compatibility with
an existing signature scheme such as Schnorr signatures.
Surprisingly, today's best non-pairing-based threshold signature
constructions that produce Schnorr signatures with unlimited
concurrency~\cite{Stinson:2001:PSD,distributed-keygen-revisited} require at
least three rounds of
communication during signing operations, whereas constructions with fewer
network rounds~\cite{distributed-keygen-revisited} must limit signing
concurrency to protect against a forgery attack~\cite{Drijvers2019OnTS}.

In this work, we present FROST,
a Flexible
Round-Optimized Schnorr Threshold
signature scheme\footnote{Signatures generated using the FROST protocol can also be
referred to as ''FROSTy signatures''.} that addresses the need for efficient
threshold signing operations while improving upon the state of the art to
ensure strong security properties \emph{without} limiting the parallelism of
signing operations.
FROST can be used as either a two-round protocol where signers send
and receive two messages in total, or optimized to a (non-broadcast)
single-round signing protocol with a pre-processing stage.
FROST achieves improved efficiency in the optimistic
case that no participant misbehaves.
However, in the case where a misbehaving participant contributes malformed
values during the protocol, honest parties can identify and exclude the
misbehaving participant, and re-run the protocol.

The flexible design of FROST lends itself to supporting a number of practical
use cases for threshold signing. Because the preprocessing round can be
performed separately from the signing round, signing operations can be performed
\emph{asynchronously};
once the preprocessing round is complete, signers only need to receive and eventually reply with a single message to create a signature.
Further, while some threshold schemes in the literature require all
participants to be active during signing
operations~\cite{distributed-keygen-revisited, fast-threshold-ecdsa}, and refer
to the threshold
property of the protocol as merely a security property, FROST allows any
threshold number of participants to produce valid signatures. Consequently,
FROST can support use cases where a subset of participants (or participating
devices) can remain offline, a property that is often desirable for
security in practice.

\medskip\noindent
\textbf{Contributions.} In this chapter, we present the
following contributions.
\begin{itemize}
  \item We review related threshold signature schemes and present a
    detailed analysis of their performance and designs.
  \item We present FROST, a Flexible Round-Optimized Schnorr Threshold signature
    scheme.
    FROST improves upon
    the state of the art for Schnorr threshold signatures by defining a
    signing protocol that can be optimized to a (non-broadcast) single-round
    operation with a preprocessing stage.
    Unlike many prior Schnorr threshold schemes, FROST remains secure against
    known forgery attacks without limiting concurrency of signing operations.
  \item We present the first proof of security and correctness for
    an interactive two-round variant of FROST, building upon proofs of security for prior
    related threshold schemes.
    We then demonstrate how this proof extends to FROST in the single-rounnd
    setting.
  \item We review follow-up proofs of security that demonstrate unforgeability for FROST directly.
\end{itemize}

